\chapter{Binary Search Trees}

You have already met a few ways to store a bag of numbers so you can find things in
them. An unsorted array lets you add a new item instantly but forces you to scan the
whole thing to find anything. A sorted array lets you find things fast with binary
search, but inserting a new item means shoving half the array over to make room. Each
structure is good at one thing and clumsy at the other.

In this chapter we build a structure that tries to be good at \emph{both}: fast lookup
\emph{and} fast insertion and deletion, all at once. It is called the \emph{binary search
tree}, or BST. The idea is to take the spirit of binary search -- ``compare, then go left
or right'' -- and freeze it into the actual shape of the data. Instead of repeatedly
cutting a sorted array in half, we build a branching structure whose every fork already
encodes a comparison.

By the end you will know how a BST is shaped, how to search it, how to find the smallest
and largest elements, how to walk from any element to the ``next'' one, and how to insert
and delete while keeping the shape valid. Then we will do something that surprises almost
everyone the first time they see it: we will show that building a BST by inserting items
in random order is, secretly, the same process as quicksort. That single observation will
hand us the expected height of a random BST almost for free.

\section{The shape of the idea}

\subsection{From sorted arrays to branching decisions}

Picture binary search on the sorted array $[2,5,8,12,15,18,21]$. To find $15$, you look at
the middle element $12$, decide $15$ is bigger, and throw away the left half. Now you look
at the middle of what remains, $18$, decide $15$ is smaller, and throw away the right half.
One more comparison and you are done. Every step is a single comparison that sends you
either left (toward smaller things) or right (toward larger things).

A binary search tree is what you get if you write that decision structure down once and
reuse it. Each element becomes a little box called a \emph{node}. Every node holds a key
and has at most two children: a \emph{left child} pointing to a subtree of smaller keys
and a \emph{right child} pointing to a subtree of larger keys. Searching is then just
walking down from the top, going left or right at each node exactly as binary search
would.

\begin{intuitionbox}
A BST is binary search turned into furniture. In a sorted array, the ``go left or right''
decision is recomputed every time from indices. In a BST, that decision is built into the
links: standing at a node, ``smaller'' literally \emph{is} the left pointer and ``larger''
literally \emph{is} the right pointer. You do not search a shape; the shape \emph{is} the
search.
\end{intuitionbox}

Here is a tree holding those same seven keys. Read it top to bottom.

\begin{center}
\begin{tikzpicture}[
  every node/.style={circle,draw=structurepurple,thick,minimum size=8mm,inner sep=1pt,font=\small},
  level distance=11mm,
  level 1/.style={sibling distance=34mm},
  level 2/.style={sibling distance=17mm},
  edge from parent/.style={draw=gray!60,thick}
]
\node {12}
  child {node {5}
    child {node {2}}
    child {node {8}}
  }
  child {node {18}
    child {node {15}}
    child {node {21}}
  };
\end{tikzpicture}
\end{center}

Notice what is true at \emph{every} node. Look at $12$: everything in its left subtree
($5,2,8$) is less than $12$, and everything in its right subtree ($18,15,21$) is greater.
Look at $5$: its left subtree ($2$) is smaller, its right ($8$) is larger. The same rule
holds locally everywhere. That single rule is the whole definition.

\begin{formalbox}
\begin{definition}[Binary search tree property]
\label{ch12:bst-property}
A \emph{binary search tree} is a binary tree in which each node $x$ stores a key
$x.key$ and pointers $x.left$, $x.right$, $x.parent$. The tree satisfies the
\emph{binary-search-tree property}: for every node $x$,
\begin{itemize}
  \item if $y$ is a node in the \textbf{left} subtree of $x$, then $y.key \le x.key$; and
  \item if $y$ is a node in the \textbf{right} subtree of $x$, then $y.key \ge x.key$.
\end{itemize}
The topmost node is the \emph{root}; a node with no children is a \emph{leaf}. A missing
child is represented by the special value $\textsc{nil}$.
\end{definition}
\end{formalbox}

\begin{trapbox}
The BST property is about \emph{entire subtrees}, rather than merely immediate children. A common
beginner error is to check only that a node's left child is smaller and its right child is
larger, and declare victory. That is not enough. In the tree below, every parent-child pair
looks locally fine, yet $9$ sits in the left subtree of $8$ where only keys $\le 8$ are
allowed:
\begin{center}
\begin{tikzpicture}[every node/.style={circle,draw=trapred,thick,minimum size=8mm,inner sep=1pt,font=\small},level distance=10mm,level 1/.style={sibling distance=24mm},level 2/.style={sibling distance=14mm}]
\node {8}
  child {node {4} child {node {1}} child {node {9}}}
  child {node {12}};
\end{tikzpicture}
\end{center}
The local check at $4$ passes ($9 > 4$, fine for a right child). But $9 > 8$, so $9$ must
not be anywhere left of $8$. This is \emph{not} a valid BST.
\end{trapbox}

\subsection{The in-order walk}

The BST property pays off immediately. Suppose you visit the whole tree in a particular
rhythm: for each node, first fully visit its \emph{left} subtree, then print the node's own
key, then fully visit its \emph{right} subtree. This rhythm is called an \emph{in-order
traversal}.

\begin{algorithm}[H]
\caption{\textsc{In-Order-Walk}($x$)}
\KwIn{A node $x$ (possibly \textsc{nil}) of a BST.}
\KwOut{Prints the keys of $x$'s subtree in sorted order.}
\If{$x \ne \textsc{nil}$}{
  \textsc{In-Order-Walk}($x.left$)\;
  print $x.key$\;
  \textsc{In-Order-Walk}($x.right$)\;
}
\end{algorithm}

Run it on the seven-key tree above. We start at $12$, but before printing $12$ we must
finish its left subtree rooted at $5$. Before printing $5$ we must finish its left subtree,
which is just $2$. So $2$ prints first. Then back up: $5$ prints, then $5$'s right child
$8$. Now the left subtree of $12$ is done, so $12$ prints. Then we descend into $12$'s right
subtree at $18$: first its left child $15$, then $18$, then $21$. The output is
\[
2,\;5,\;8,\;12,\;15,\;18,\;21,
\]
which is \emph{exactly} sorted order. This is not a coincidence; it is the BST property
speaking out loud.

Figure~\ref{ch12:fig-inorder} shows the visiting order as small superscripts: each node is
printed only after its entire left subtree, so the keys come out smallest-first.

\begin{figure}[H]
\centering
\begin{tikzpicture}[
  every node/.style={circle,draw=structurepurple,thick,minimum size=9mm,inner sep=1pt,font=\small},
  level distance=12mm,
  level 1/.style={sibling distance=34mm},
  level 2/.style={sibling distance=17mm},
  edge from parent/.style={draw=gray!60,thick}
]
\node {$12^{4}$}
  child {node {$5^{2}$}
    child {node {$2^{1}$}}
    child {node {$8^{3}$}}
  }
  child {node {$18^{6}$}
    child {node {$15^{5}$}}
    child {node {$21^{7}$}}
  };
\end{tikzpicture}
\caption{The in-order walk of the seven-key tree. Superscripts give the print order:
left subtree, then the node, then the right subtree. Reading the superscripts $1,2,\dots,7$
visits the keys $2,5,8,12,15,18,21$ in sorted order.}
\label{ch12:fig-inorder}
\end{figure}

\begin{focusbox}
An in-order walk of a BST prints the keys in sorted order. So a BST is, among other things,
a way of \emph{storing a sorted sequence} that you can also splice into and delete from
cheaply. Keep the picture in mind: ``left subtree, me, right subtree'' = ``everything smaller,
me, everything larger'' = sorted.
\end{focusbox}

\begin{lemma}[Correctness of the in-order walk]
\label{ch12:inorder-correct}
\textsc{In-Order-Walk}$(x)$ prints the keys of the subtree rooted at $x$ in nondecreasing
order.
\end{lemma}

\begin{proof}
We argue by induction on the size $n$ of the subtree rooted at $x$. If $n = 0$ then $x =
\textsc{nil}$ and nothing is printed; an empty sequence is vacuously sorted, so the base case
holds. Now suppose the claim holds for all subtrees of size less than $n$, and let $x$ have
a subtree of size $n \ge 1$. Its left and right subtrees each have fewer than $n$ nodes, so
by the induction hypothesis the call \textsc{In-Order-Walk}$(x.left)$ prints the left
subtree's keys in sorted order, and likewise on the right. By the BST property
(Definition~\ref{ch12:bst-property}), every left-subtree key is $\le x.key$ and every
right-subtree key is $\ge x.key$. The walk prints (left keys, sorted) then $x.key$ then
(right keys, sorted), and these three blocks are already in the correct relative order. The
concatenation is therefore sorted.
\end{proof}

Since the walk touches each of the $n$ nodes exactly once and does a constant amount of
work per node, it runs in $\Theta(n)$ time.

\begin{walkbox}
You have just seen that a BST can be flattened into a sorted list in linear time. So here is
a question to chew on while you walk: \emph{given} a sorted list, how many different BSTs
produce it as their in-order walk? The seven-key tree we drew was nicely balanced, but is it
the only tree whose in-order walk is $2,5,8,12,15,18,21$? Try to build a wildly different one
by hand. What does that tell you about how much the \emph{shape} of a BST can vary while the
\emph{contents} stay fixed?
\end{walkbox}

\section{Searching, minimum, maximum}

\subsection{Searching for a key}

Searching is the operation a BST is named for, and it is exactly the binary-search rhythm.
Start at the root. If the node's key equals what you want, stop. If your target is smaller,
go to the left child; if larger, go to the right. Repeat until you find the key or fall off
the bottom of the tree into a \textsc{nil}.

\begin{algorithm}[H]
\caption{\textsc{Tree-Search}($x$, $k$)}
\KwIn{Root $x$ of a (sub)tree, search key $k$.}
\KwOut{A node with key $k$, or \textsc{nil} if absent.}
\While{$x \ne \textsc{nil}$ \textbf{and} $k \ne x.key$}{
  \eIf{$k < x.key$}{$x \gets x.left$}{$x \gets x.right$}
}
\Return $x$\;
\end{algorithm}

Let us trace a search for $k = 15$ in the seven-key tree. We start at $12$. Since $15 > 12$
we go right to $18$. Since $15 < 18$ we go left to $15$. Now $k = x.key$, so we stop and
return that node. Three comparisons, one per level. A search for $k = 7$ would go: $12$ (go
left, $7<12$) $\to 5$ (go right, $7>5$) $\to 8$ (go left, $7<8$) $\to \textsc{nil}$, and we
return \textsc{nil} -- $7$ is not in the tree. Figure~\ref{ch12:fig-search} highlights both
paths: the successful hunt for $15$ ending on a node, and the failed hunt for $7$ running off
the bottom into a missing child.

\begin{figure}[H]
\centering
\begin{tikzpicture}[
  every node/.style={circle,draw=structurepurple,thick,minimum size=8mm,inner sep=1pt,font=\small},
  level distance=12mm,
  level 1/.style={sibling distance=34mm},
  level 2/.style={sibling distance=17mm},
  edge from parent/.style={draw=gray!55,thick}
]
\node (s12) {12}
  child {node (s5) {5}
    child {node {2}}
    child {node (s8) {8}}
  }
  child {node (s18) {18}
    child {node (s15) {15}}
    child {node {21}}
  };
% successful path to 15 (green)
\draw[-{Stealth[length=2.2mm]},targetgreen,very thick] (s12) to[bend left=8] (s18);
\draw[-{Stealth[length=2.2mm]},targetgreen,very thick] (s18) to[bend left=8] (s15);
% failed path to 7 (red), down from 8 into a missing left child
\draw[-{Stealth[length=2.2mm]},trapred,very thick] (s12) to[bend right=8] (s5);
\draw[-{Stealth[length=2.2mm]},trapred,very thick] (s5) to[bend right=8] (s8);
\draw[-{Stealth[length=2.2mm]},trapred,very thick] (s8) -- ++(-0.7,-1.05)
  node[rectangle,minimum size=0pt,draw=none,fill=none,below,font=\scriptsize\itshape,trapred] {\textsc{nil}};
\node[rectangle,minimum size=0pt,draw=none,fill=none,font=\footnotesize,targetgreen] at (2.8,-3.1) {find $15$: hit};
\node[rectangle,minimum size=0pt,draw=none,fill=none,font=\footnotesize,trapred] at (-2.8,-3.1) {find $7$: miss};
\end{tikzpicture}
\caption{Two searches in the seven-key tree. Green: searching for $15$ goes right then left
and lands on a node. Red: searching for $7$ goes left, right, left, then falls off the bottom
into \textsc{nil}, so $7$ is absent.}
\label{ch12:fig-search}
\end{figure}

\begin{intuitionbox}
A search never wanders. At each node it either finishes or commits to exactly one child and
never looks back. So a search traces a single \emph{path} from the root downward. The number
of comparisons is at most the length of that path, which is at most the \emph{height} of the
tree. This is why height is the number that controls everything in this chapter.
\end{intuitionbox}

\subsection{Minimum and maximum}

Where is the smallest key? The BST property says everything smaller than a node lives in its
left subtree. So to find the minimum, keep following left children until you cannot anymore.
The last node you reach has no left child, which means nothing in the tree is smaller. The
maximum is the mirror image: follow right children to the end.

\begin{algorithm}[H]
\caption{\textsc{Tree-Minimum}($x$)}
\KwIn{Root $x \ne \textsc{nil}$ of a (sub)tree.}
\KwOut{The node holding the smallest key.}
\While{$x.left \ne \textsc{nil}$}{$x \gets x.left$}
\Return $x$\;
\end{algorithm}

In our tree the minimum walk is $12 \to 5 \to 2$, and $2$ has no left child, so $2$ is the
minimum. The maximum walk is $12 \to 18 \to 21$, giving $21$. Both touch one node per level,
so both run in $O(h)$ time where $h$ is the height.

\section{The successor: finding the next key}

Suppose you are standing at a node and you want the node with the \emph{next-largest} key --
the one that would come immediately after it in the sorted in-order list. This is the
\emph{successor}, and it is the trickiest of the basic queries, so let us build it carefully.

There are two cases, and the picture makes them obvious. Think about where the ``next''
element hides relative to the in-order walk.

\textbf{Case 1: the node has a right subtree.} Then everything in that right subtree is
larger than our node, and the successor is the \emph{smallest} thing in there. So the answer
is \textsc{Tree-Minimum} of the right child. For example, the successor of $12$ is the
minimum of the right subtree rooted at $18$, which is $15$. Check the sorted list:
$\ldots,12,15,\ldots$ -- correct.

\textbf{Case 2: the node has no right subtree.} Then there is nothing larger ``below'' it, so
the next-largest key must be ``above'' it. We climb upward. As long as we are coming up from
a \emph{right} child, our parent is smaller than us, so it cannot be the successor; keep
climbing. The first time we come up from a \emph{left} child, that parent is larger than us
and is the successor. (If we climb all the way off the top, our node was the maximum and has
no successor.)

Figure~\ref{ch12:fig-successor} shows both cases on the seven-key tree. On the left, the
successor of $12$ is found by descending to the minimum of its right subtree. On the right,
the successor of $8$ (which has no right child) is found by climbing: up through $5$ as a
right child, then stopping at $12$, the first ancestor reached from a left child.

\begin{figure}[H]
\centering
\begin{tikzpicture}[
  every node/.style={circle,draw=structurepurple,thick,minimum size=8mm,inner sep=1pt,font=\small},
  level distance=11mm,
  level 1/.style={sibling distance=22mm},
  level 2/.style={sibling distance=12mm},
  edge from parent/.style={draw=gray!55,thick}
]
% Case 1: successor of 12 = min of right subtree
\begin{scope}
\node[fill=sourceblue!18] (a12) {12}
  child {node {5} child {node {2}} child {node {8}}}
  child {node (a18) {18} child {node[fill=targetgreen!30] (a15) {15}} child {node {21}}};
\node[draw=none,fill=none,font=\footnotesize\itshape] at (0,0.9) {start (12)};
\draw[-{Stealth[length=2mm]},targetgreen,thick] (a12) to[bend left=12] (a18);
\draw[-{Stealth[length=2mm]},targetgreen,thick] (a18) to[bend left=12] (a15);
\node[draw=none,fill=none,font=\footnotesize\itshape,targetgreen] at (1.0,-3.5) {successor};
\node[draw=none,fill=none,font=\footnotesize] at (0,-3.9) {Case 1: go right, then all the way left};
\end{scope}
% Case 2: successor of 8 = climb to first left-ancestor
\begin{scope}[xshift=62mm]
\node (b12) {12}
  child {node (b5) {5} child {node {2}} child {node[fill=sourceblue!18] (b8) {8}}}
  child {node {18} child {node {15}} child {node {21}}};
\draw[-{Stealth[length=2mm]},quotientorange,thick] (b8) to[bend left=18] (b5);
\draw[-{Stealth[length=2mm]},quotientorange,thick] (b5) to[bend left=12] (b12);
\node[draw=none,fill=none,font=\footnotesize\itshape] at (1.6,-2.6) {start (8)};
\node[draw=none,fill=none,font=\footnotesize\itshape,quotientorange] at (-0.2,0.95) {successor};
\node[draw=none,fill=none,font=\footnotesize] at (0,-3.9) {Case 2: climb until you rise from a left child};
\end{scope}
\end{tikzpicture}
\caption{The two cases of \textsc{Tree-Successor}. Left: when a right subtree exists, the
successor is its minimum. Right: with no right subtree, the successor is the lowest ancestor
whose left subtree contains the node.}
\label{ch12:fig-successor}
\end{figure}

\begin{algorithm}[H]
\caption{\textsc{Tree-Successor}($x$)}
\KwIn{A node $x$ in a BST.}
\KwOut{The node with the next-larger key, or \textsc{nil} if $x$ is the maximum.}
\If{$x.right \ne \textsc{nil}$}{
  \Return \textsc{Tree-Minimum}($x.right$)\;
}
$y \gets x.parent$\;
\While{$y \ne \textsc{nil}$ \textbf{and} $x = y.right$}{
  $x \gets y$\;
  $y \gets y.parent$\;
}
\Return $y$\;
\end{algorithm}

\begin{labbox}
Trace \textsc{Tree-Successor} on node $8$ in the seven-key tree. Node $8$ has no right child,
so we are in Case 2. Set $y = 5$ (the parent of $8$). Is $8$ a right child of $5$? Yes, so we
climb: now $x = 5$, $y = 12$. Is $5$ a right child of $12$? No, $5$ is the \emph{left} child of
$12$. The loop stops and we return $y = 12$. And indeed, in $2,5,8,12,\ldots$ the key right
after $8$ is $12$. Now trace it yourself on node $21$ (the maximum): you should climb all the
way up, fall off the top, and return \textsc{nil}.
\end{labbox}

Both cases follow a single path -- either straight down the right subtree or straight up the
ancestor chain -- so \textsc{Tree-Successor} runs in $O(h)$ time. The mirror algorithm,
\textsc{Tree-Predecessor}, finds the next-\emph{smaller} key by swapping ``left'' and
``right.''

\begin{aestheticbox}
Successor and predecessor are what make a BST behave like a navigable, ordered collection
rather than a mere lookup table. They are why the ordered maps and sets in standard libraries
(C++ \lstinline|std::map|, Java's \lstinline|TreeMap|) can answer ``give me the smallest key
$\ge x$'' or iterate keys in sorted order. A hash table, fast as it is, cannot do any of that:
it has no notion of ``next.'' When you need range queries or ordered iteration, you reach for
a balanced BST.
\end{aestheticbox}

\section{Insertion}

To insert a new key, we want to find the spot where a \emph{search} for that key would have
fallen off the tree into a \textsc{nil}, and hang the new node there. That placement
automatically preserves the BST property, because we walked down using the very comparisons
the property demands.

We keep two pointers as we descend: $x$ races ahead looking for the empty slot, and a
trailing pointer $y$ remembers $x$'s parent so we know what to attach the new node to.

\begin{algorithm}[H]
\caption{\textsc{Tree-Insert}($T$, $z$)}
\KwIn{Tree $T$ and a new node $z$ with $z.key$ set, $z.left = z.right = \textsc{nil}$.}
\KwOut{$T$ with $z$ inserted.}
$y \gets \textsc{nil}$\;
$x \gets T.root$\;
\While{$x \ne \textsc{nil}$}{
  $y \gets x$\;
  \eIf{$z.key < x.key$}{$x \gets x.left$}{$x \gets x.right$}
}
$z.parent \gets y$\;
\uIf{$y = \textsc{nil}$}{$T.root \gets z$ \tcp*{tree was empty}}
\uElseIf{$z.key < y.key$}{$y.left \gets z$}
\Else{$y.right \gets z$}
\end{algorithm}

\begin{example}[Building a tree by hand]
\label{ch12:insert-trace}
Start with an empty tree and insert keys $12, 5, 18, 2, 8, 15, 21$ in order.
\begin{itemize}
  \item Insert $12$: tree is empty, so $12$ becomes the root.
  \item Insert $5$: at root $12$, $5 < 12$, go left into \textsc{nil}. Attach $5$ as the left
        child of $12$.
  \item Insert $18$: at $12$, $18 > 12$, go right into \textsc{nil}. Attach $18$ as right child.
  \item Insert $2$: $12 \to$ left to $5 \to$ ($2<5$) left into \textsc{nil}. Attach under $5$.
  \item Insert $8$: $12 \to$ left to $5 \to$ ($8>5$) right into \textsc{nil}. Attach under $5$.
  \item Insert $15$: $12 \to$ right to $18 \to$ ($15<18$) left into \textsc{nil}. Attach under $18$.
  \item Insert $21$: $12 \to$ right to $18 \to$ ($21>18$) right into \textsc{nil}. Attach under $18$.
\end{itemize}
The result is exactly the balanced seven-key tree from the start of the chapter. Notice that
the root is whatever you inserted \emph{first}. We will lean on that fact heavily later.
\end{example}

\begin{figure}[h]
\centering
\begin{tikzpicture}[font=\small]
  % Panel A: Trailing pointer descent
  \begin{scope}
    \node[font=\footnotesize\scshape, text=gray!70!black] at (0, 1.2) {(a) Descent path for $z=13$};
    \node[circle, draw=structurepurple, fill=white, thick, minimum size=7.5mm] (n12) at (0,0) {12};
    \node[circle, draw=structurepurple, fill=white, thick, minimum size=7.5mm] (n5) at (-1.5,-1.1) {5};
    \node[circle, draw=structurepurple, fill=white, thick, minimum size=7.5mm] (n18) at (1.5,-1.1) {18};
    \node[circle, draw=structurepurple, fill=white, thick, minimum size=7.5mm] (n2) at (-2.1,-2.2) {2};
    \node[circle, draw=structurepurple, fill=white, thick, minimum size=7.5mm] (n8) at (-0.9,-2.2) {8};
    \node[circle, draw=structurepurple, fill=white, thick, minimum size=7.5mm] (n15) at (0.9,-2.2) {15};
    \node[circle, draw=structurepurple, fill=white, thick, minimum size=7.5mm] (n21) at (2.1,-2.2) {21};
    
    \draw[thick, draw=gray!60] (n12) -- (n5);
    \draw[thick, draw=gray!60] (n12) -- (n18);
    \draw[thick, draw=gray!60] (n5) -- (n2);
    \draw[thick, draw=gray!60] (n5) -- (n8);
    \draw[thick, draw=gray!60] (n18) -- (n15);
    \draw[thick, draw=gray!60] (n18) -- (n21);

    % Animated search path
    \draw[-{Latex[length=2.5mm]}, very thick, trapred] (n12) to[bend left=15] node[right, font=\tiny, text=trapred] {$13 > 12$} (n18);
    \draw[-{Latex[length=2.5mm]}, very thick, trapred] (n18) to[bend right=15] node[left, font=\tiny, text=trapred] {$13 < 18$} (n15);
    \draw[-{Latex[length=2.5mm]}, very thick, trapred] (n15) -- node[left, font=\tiny, text=trapred] {$13 < 15$} (0.3,-3.1);
    \node[draw=gray!50, dashed, rounded corners=2pt, font=\tiny, inner sep=2pt] at (0.3,-3.3) {\textsc{nil}};
    
    \node[font=\scriptsize, text=sourceblue!80!black, align=center] at (0, -4.2) {Trailing pointer $y$ stops at $15$};
  \end{scope}

  % Panel B: Grafting the new leaf
  \begin{scope}[xshift=7.2cm]
    \node[font=\footnotesize\scshape, text=gray!70!black] at (0, 1.2) {(b) Attaching new leaf $z$};
    \node[circle, draw=structurepurple, fill=white, thick, minimum size=7.5mm] (b12) at (0,0) {12};
    \node[circle, draw=structurepurple, fill=white, thick, minimum size=7.5mm] (b5) at (-1.5,-1.1) {5};
    \node[circle, draw=structurepurple, fill=white, thick, minimum size=7.5mm] (b18) at (1.5,-1.1) {18};
    \node[circle, draw=structurepurple, fill=white, thick, minimum size=7.5mm] (b2) at (-2.1,-2.2) {2};
    \node[circle, draw=structurepurple, fill=white, thick, minimum size=7.5mm] (b8) at (-0.9,-2.2) {8};
    \node[circle, draw=sourceblue, fill=sourceblue!15, thick, minimum size=7.5mm] (b15) at (0.9,-2.2) {15};
    \node[circle, draw=structurepurple, fill=white, thick, minimum size=7.5mm] (b21) at (2.1,-2.2) {21};
    \node[circle, draw=targetgreen, fill=targetgreen!20, thick, minimum size=7.5mm] (b13) at (0.3,-3.3) {13};
    
    \draw[thick, draw=gray!60] (b12) -- (b5);
    \draw[thick, draw=gray!60] (b12) -- (b18);
    \draw[thick, draw=gray!60] (b5) -- (b2);
    \draw[thick, draw=gray!60] (b5) -- (b8);
    \draw[thick, draw=gray!60] (b18) -- (b15);
    \draw[thick, draw=gray!60] (b18) -- (b21);
    \draw[very thick, draw=targetgreen] (b15) -- (b13);
    
    \node[font=\scriptsize\bfseries, text=sourceblue!90!black, above right=1pt of b15] {$y$};
    \node[font=\scriptsize\bfseries, text=targetgreen!90!black, left=2pt of b13] {$z$};
    \node[font=\scriptsize, text=targetgreen!80!black, align=center] at (0, -4.2) {BST property: $12 < 13 < 15 < 18$\\[1pt]\scriptsize\color{sourceblue!80!black}$y.\mathit{left} \gets z$, $z.p \gets y$};
  \end{scope}
\end{tikzpicture}
\caption{The mechanics of \textsc{Tree-Insert}($T, z$) with $z.key = 13$. In (a), pointer $x$ walks down the tree guided by key comparisons, while trailing pointer $y$ lags one step behind. In (b), when $x$ falls off the tree at \textsc{nil}, $z$ is spliced in as $y$'s left child.}
\label{ch12:fig-insert-walk}
\end{figure}

Insertion walks one root-to-leaf path and then does constant work to attach the node, so it
costs $O(h)$.

\begin{trapbox}
The order you insert keys determines the shape you get. Insert $2,5,8,12,15,18,21$ -- already
sorted -- and every key is larger than the last, so each one becomes the right child of the
previous one. You get a ``tree'' that is really a straight line slanting to the right:
\begin{center}
\begin{tikzpicture}[every node/.style={circle,draw=trapred,thick,minimum size=7mm,inner sep=1pt,font=\footnotesize},level distance=8mm,level 1/.style={sibling distance=10mm}]
\node {2} child[missing] child {node {5} child[missing] child {node {8} child[missing] child {node {12} child[missing] child {node {15} child[missing] child {node {18} child[missing] child {node {21}}}}}}};
\end{tikzpicture}
\end{center}
This is a valid BST, but its height is $n-1$, so searching it is no better than scanning a
linked list: $\Theta(n)$ per operation. Same keys, catastrophically worse shape. The whole
next chapter (red-black trees) exists to prevent exactly this.
\end{trapbox}

\begin{intuitionbox}
\textbf{How many BST shapes can $n$ keys form? (The Catalan numbers.)}
If the insertion sequence dictates the tree's geometry, how vast is the universe of possible shapes?
Suppose we have $n$ distinct keys, which we can label $1, 2, \dots, n$ in increasing order. Any key $i \in \{1, \dots, n\}$ could end up as the root. Because every key smaller than $i$ must go into the left subtree and every key larger than $i$ into the right subtree, choosing key $i$ as the root leaves exactly $i-1$ keys on the left and $n-i$ keys on the right.

The shapes of the left and right subtrees can be chosen independently. If $C_n$ denotes the number of distinct binary tree structures with $n$ nodes, we obtain the fundamental convolution recurrence:
\[
  C_n = \sum_{i=1}^n C_{i-1} C_{n-i}, \quad \text{with } C_0 = 1.
\]
For small values: $C_1 = 1$, $C_2 = C_0 C_1 + C_1 C_0 = 2$, $C_3 = C_0 C_2 + C_1 C_1 + C_2 C_0 = 2 + 1 + 2 = 5$.
These combinatorial quantities are the \emph{Catalan numbers}, with the exact closed-form formula:
\[
  C_n = \frac{1}{n+1}\binom{2n}{n}.
\]
As $n$ grows, Stirling's approximation reveals $C_n \sim \frac{4^n}{n^{3/2}\sqrt{\pi}}$. For $n = 10$, there are $16{,}796$ distinct tree shapes; for $n = 20$, over $6.5 \times 10^9$. Among this combinatorial explosion of shapes, only a vanishingly tiny fraction are perfectly balanced, while $2^{n-1}$ of them are degenerate linear sticks. Fortunately, when keys arrive in random order, the probability concentrates heavily on trees with height $O(\ln n)$---as we will prove in Section~12.7.
\end{intuitionbox}

\section{Deletion}

Deletion is the fussiest BST operation because removing a node can leave a hole in the middle
of the structure, and we have to fill that hole without breaking the BST property. The good
news is there are only three cases, and the third reduces to the others.

Suppose we want to delete node $z$.

\textbf{Case A: $z$ has no children.} It is a leaf. Just detach it -- replace the pointer
from its parent with \textsc{nil}. Nothing else can break.

\textbf{Case B: $z$ has exactly one child.} Lift that child up to take $z$'s place,
reconnecting it to $z$'s parent. Because the single child subtree was already entirely on the
correct side of $z$'s ancestors, sliding it up one level keeps everything ordered.

\textbf{Case C: $z$ has two children.} This is the interesting case. We cannot simply promote a single
child, because $z$ has two branches requiring a parent slot. Instead we find a key that can sit in
$z$'s position without disturbing the order: the \emph{successor} of $z$, call it $y$. The
successor is the smallest key in $z$'s right subtree, so it is larger than everything in
$z$'s left subtree and smaller than everything else in $z$'s right subtree -- precisely the
property $z$'s key had. We splice $y$ out of its current spot and let it stand in for $z$.

\begin{intuitionbox}
Why the successor specifically? Because we need a replacement that is $\ge$ all of $z$'s left
subtree and $\le$ all the rest of $z$'s right subtree. The successor -- the next key up in
sorted order -- is the unique tightest fit on the right side. Moreover, the successor of
a two-child node lives in its right subtree and has \emph{no left child} (it is the minimum
there), so plucking it out is easy: it falls under Case A or B.
\end{intuitionbox}

To keep the code clean, we use a small helper, \textsc{Transplant}, that replaces the subtree
rooted at $u$ with the subtree rooted at $v$ by fixing up exactly one parent pointer.

\begin{algorithm}[H]
\caption{\textsc{Transplant}($T$, $u$, $v$)}
\KwIn{Tree $T$; node $u$ to be replaced; node $v$ (possibly \textsc{nil}) to put in its place.}
\uIf{$u.parent = \textsc{nil}$}{$T.root \gets v$}
\uElseIf{$u = u.parent.left$}{$u.parent.left \gets v$}
\Else{$u.parent.right \gets v$}
\If{$v \ne \textsc{nil}$}{$v.parent \gets u.parent$}
\end{algorithm}

\begin{figure}[h]
\centering
\begin{tikzpicture}[font=\small]
  % Before transplant
  \begin{scope}
    \node[font=\footnotesize\scshape, text=gray!70!black] at (1.5, 2.2) {Before: \textsc{Transplant}($T, u, v$)};
    \node[circle, draw=gray!70, fill=gray!10, thick, minimum size=8mm] (p) at (1.5, 1.2) {$u.p$};
    \node[circle, draw=trapred, fill=trapred!15, thick, minimum size=8mm] (u) at (0.5, 0) {$u$};
    \node[circle, draw=gray!40, fill=white, dashed, minimum size=7mm] (sib) at (2.5, 0) {sib};
    
    \draw[thick, draw=gray!60] (p) -- (sib);
    \draw[very thick, draw=trapred] (p) -- node[left, font=\tiny, text=trapred] {$u.p.left$} (u);
    
    % Subtrees of u
    \draw[thick, draw=trapred!70, fill=trapred!5] (0.1, -0.6) -- (-0.4, -1.5) -- (0.6, -1.5) -- cycle;
    \node[font=\tiny, text=trapred!80!black] at (0.1, -1.1) {$T_u$};
    
    % Subtree v standing elsewhere
    \node[circle, draw=targetgreen, fill=targetgreen!20, thick, minimum size=8mm] (v) at (3.5, 0) {$v$};
    \draw[thick, draw=targetgreen!70, fill=targetgreen!5] (3.5, -0.6) -- (3.0, -1.5) -- (4.0, -1.5) -- cycle;
    \node[font=\tiny, text=targetgreen!80!black] at (3.5, -1.1) {$T_v$};
  \end{scope}

  % Transition arrow
  \draw[-{Latex[length=3mm]}, very thick, quotientorange] (4.6, 0.3) -- node[above, font=\scriptsize, text=quotientorange] {Rewire pointers} (5.8, 0.3);

  % After transplant
  \begin{scope}[xshift=6.8cm]
    \node[font=\footnotesize\scshape, text=gray!70!black] at (1.5, 2.2) {After: $v$ takes $u$'s slot under $u.p$};
    \node[circle, draw=gray!70, fill=gray!10, thick, minimum size=8mm] (p2) at (1.5, 1.2) {$u.p$};
    \node[circle, draw=targetgreen, fill=targetgreen!20, thick, minimum size=8mm] (v2) at (0.5, 0) {$v$};
    \node[circle, draw=gray!40, fill=white, dashed, minimum size=7mm] (sib2) at (2.5, 0) {sib};
    
    \draw[thick, draw=gray!60] (p2) -- (sib2);
    \draw[very thick, draw=targetgreen] (p2) -- node[left, font=\tiny, text=targetgreen!80!black] {$u.p.left \gets v$} (v2);
    \draw[-{Latex[length=2mm]}, very thick, targetgreen] (v2) to[bend left=25] node[right, font=\tiny, text=targetgreen!80!black] {$v.p \gets u.p$} (p2);
    
    % Subtree v attached
    \draw[thick, draw=targetgreen!70, fill=targetgreen!5] (0.5, -0.6) -- (0.0, -1.5) -- (1.0, -1.5) -- cycle;
    \node[font=\tiny, text=targetgreen!80!black] at (0.5, -1.1) {$T_v$};
    
    \node[font=\scriptsize, text=gray!70!black] at (1.5, -2.0) {$u$ is detached; $v$ inherits $u$'s position};
  \end{scope}
\end{tikzpicture}
\caption{The \textsc{Transplant}($T, u, v$) subroutine replaces the subtree rooted at $u$ with the subtree rooted at $v$. It updates $u$'s parent to point to $v$ and sets $v$'s parent pointer to $u.p$, leaving $u$'s former children undisturbed for caller-specific handling.}
\label{ch12:fig-transplant}
\end{figure}

\begin{algorithm}[H]
\caption{\textsc{Tree-Delete}($T$, $z$)}
\KwIn{Tree $T$ and a node $z$ to remove.}
\uIf{$z.left = \textsc{nil}$}{
  \textsc{Transplant}($T$, $z$, $z.right$) \tcp*{0 or 1 child: lift the right}
}
\uElseIf{$z.right = \textsc{nil}$}{
  \textsc{Transplant}($T$, $z$, $z.left$) \tcp*{exactly 1 child: lift the left}
}
\Else{
  $y \gets \textsc{Tree-Minimum}(z.right)$ \tcp*{$y$ = successor of $z$}
  \If{$y.parent \ne z$}{
    \textsc{Transplant}($T$, $y$, $y.right$)\;
    $y.right \gets z.right$\;
    $y.right.parent \gets y$\;
  }
  \textsc{Transplant}($T$, $z$, $y$)\;
  $y.left \gets z.left$\;
  $y.left.parent \gets y$\;
}
\end{algorithm}

The first two branches cover Cases A and B together (if the left child is \textsc{nil}, lifting
the right child handles both a leaf and a one-right-child node; the \emph{else if} handles a
lone left child). The final branch is Case C: we find the successor $y$, detach it from its old
location, and slot it in where $z$ was, inheriting both of $z$'s subtrees.

\begin{example}[Deleting a two-child node]
\label{ch12:delete-trace}
Delete $12$ (the root) from the balanced seven-key tree. Node $12$ has two children, so we are
in Case C. Its successor is the minimum of the right subtree rooted at $18$: follow left from
$18$ to $15$, which has no left child. So $y = 15$.

Since $y$'s parent is $18$, not $12$, we first transplant $y$ out: $15$ has no right child, so
$18$'s left pointer becomes \textsc{nil}, and we set $15$'s right child to be $z$'s old right
subtree ($18$). Then we transplant $15$ into $12$'s place at the root and give it $12$'s left
subtree ($5$). Figure~\ref{ch12:fig-delete} shows the tree just before and just after the
surgery.
\begin{figure}[H]
\centering
\begin{tikzpicture}[
  every node/.style={circle,draw=structurepurple,thick,minimum size=8mm,inner sep=1pt,font=\small},
  level distance=11mm,
  level 1/.style={sibling distance=24mm},
  level 2/.style={sibling distance=13mm},
  edge from parent/.style={draw=gray!60,thick}
]
% Before
\begin{scope}
\node[fill=trapred!22] {12}
  child {node {5} child {node {2}} child {node {8}}}
  child {node {18} child {node[fill=targetgreen!30] {15}} child {node {21}}};
\node[draw=none,fill=none,font=\footnotesize] at (0,-3.7) {before: delete $12$ (red),};
\node[draw=none,fill=none,font=\footnotesize] at (0,-4.15) {successor $15$ (green)};
\end{scope}
% After
\begin{scope}[xshift=58mm]
\node[fill=targetgreen!30] {15}
  child {node {5} child {node {2}} child {node {8}}}
  child {node {18} child[missing] child {node {21}}};
\node[draw=none,fill=none,font=\footnotesize] at (0,-3.7) {after: $15$ rises to the root,};
\node[draw=none,fill=none,font=\footnotesize] at (0,-4.15) {keeping both old subtrees};
\end{scope}
\draw[-{Stealth[length=2.5mm]},gray!70,very thick] (2.6,-1.6) -- (3.7,-1.6);
\end{tikzpicture}
\caption{Deleting the two-child node $12$. Its successor $15$ (the minimum of the right
subtree) is spliced out of its old spot and lifted into $12$'s place, inheriting both of
$12$'s subtrees. The in-order order is preserved.}
\label{ch12:fig-delete}
\end{figure}
Read the in-order walk of the result: $2,5,8,15,18,21$. Still sorted, with $12$ gone. The
property survived.
\end{example}

Like every operation in this chapter, deletion does a constant amount of pointer surgery plus
at most one \textsc{Tree-Minimum} call, so it runs in $O(h)$.

\begin{complexitybox}
Every dynamic-set operation on a BST -- search, minimum, maximum, successor, predecessor,
insert, delete -- runs in $O(h)$ time, where $h$ is the height of the tree.
\begin{center}
\begin{tabular}{lll}
\toprule
Operation & Balanced & Worst (degenerate) \\
\midrule
\textsc{Search}/\textsc{Min}/\textsc{Max} & $\Theta(\log n)$ & $\Theta(n)$ \\
\textsc{Successor}/\textsc{Predecessor}   & $\Theta(\log n)$ & $\Theta(n)$ \\
\textsc{Insert}/\textsc{Delete}           & $\Theta(\log n)$ & $\Theta(n)$ \\
\textsc{In-Order-Walk} (all keys)         & $\Theta(n)$      & $\Theta(n)$ \\
\bottomrule
\end{tabular}
\end{center}
Everything hinges on $h$. A balanced tree has $h = \Theta(\log n)$; a degenerate one has
$h = \Theta(n)$. The rest of the chapter is about which one you should \emph{expect}.
\end{complexitybox}

\section{BSTs and quicksort are the same algorithm}

\subsection{The first surprise}

Look again at Example~\ref{ch12:insert-trace}, where we built a tree by inserting
$12,5,18,2,8,15,21$. The very first key, $12$, became the root. Then every later key got
compared against $12$ on its way down: the smaller ones ($5,2,8$) drifted into the left
subtree and the larger ones ($18,15,21$) into the right subtree.

Now recall quicksort. You pick a \emph{pivot}, then partition the rest of the array into
``everything smaller than the pivot'' and ``everything larger,'' and recurse on each side.

Here is the connection. \emph{The root of a BST plays the exact role of the first pivot in
quicksort.} The first inserted key splits all the others into a smaller group and a larger group, just
like a pivot. The root of the left subtree is then the pivot for the smaller group, the root
of the right subtree is the pivot for the larger group, and so on, recursively. The set of
comparisons quicksort makes is, key for key, the set of comparisons that BST insertion makes.

\begin{center}
\begin{tikzpicture}[
  every node/.style={circle,draw=sourceblue,thick,minimum size=8mm,inner sep=1pt,font=\small},
  level distance=12mm,
  level 1/.style={sibling distance=36mm},
  level 2/.style={sibling distance=18mm},
  edge from parent/.style={draw=gray!55,thick}
]
\node {12}
  child {node {5} child {node {2}} child {node {8}}}
  child {node {18} child {node {15}} child {node {21}}};
\node[draw=none,fill=none,font=\footnotesize\itshape,structurepurple] at (0,0.9) {pivot};
\node[draw=none,fill=none,font=\footnotesize\itshape,structurepurple] at (-3.4,-0.6) {sub-pivot};
\node[draw=none,fill=none,font=\footnotesize\itshape,structurepurple] at (3.4,-0.6) {sub-pivot};
\end{tikzpicture}
\end{center}

\begin{focusbox}
Building a BST by inserting keys in some order, and running quicksort using those same keys
as pivots in that same order, perform \emph{exactly the same comparisons between the same
pairs of keys}. The BST is a recording of quicksort's recursion tree; the in-order walk of the
BST replays quicksort's final sorted output.
\end{focusbox}

\begin{historybox}
The binary search tree was independently formulated around 1960 by P.~F.~Windley, A.~D.~Booth, and A.~J.~T.~Colin to bridge the gap between sorted arrays (which allow $O(\log n)$ binary search but require $\Theta(n)$ shifts on insertion) and linked lists (which allow $O(1)$ insertion once located but require $\Theta(n)$ sequential scans). By structuring nodes as a linked hierarchy preserving left-smaller and right-larger invariants, BSTs aimed to deliver logarithmic operations dynamically. In 1962, Thomas N.~Hibbard published the classic deletion algorithm using the in-order successor. Donald Knuth (1971) later developed the dynamic programming formulation for optimal static search trees with known access probabilities.

For years, practitioners assumed that if keys arrived in random order, a BST would perpetually maintain its $O(\log n)$ average depth even as elements were repeatedly inserted and deleted. This intuition was shattered by empirical simulations and later proved by Jeffrey Culberson and J.~Ian Munro (1989): Hibbard deletion is structurally asymmetric. Replacing a deleted node with its in-order successor systematically unbalances subtrees, causing the average path length to drift away from $O(\log n)$ toward $\Theta(\sqrt{n})$ after long sequences of random updates. Furthermore, unaugmented BSTs are acutely vulnerable to natural sorted or nearly sorted input streams, which degenerate tree height to $n-1$, reducing the structure to an ordinary linked list.

Dynamic maintenance of sorted order without pre-allocating contiguous blocks. Unlike hash tables, a BST efficiently supports sorted order traversals, minimum/maximum queries, predecessors, successors, and range lookups in time proportional to tree height $h$. Yet this carries an honest trade-off: without active rebalancing mechanisms, BST worst-case time is $\Theta(n)$. Every query is hostage to insertion history, and Hibbard deletion slowly erodes tree balance over time. To guarantee $O(\log n)$ worst-case bounds regardless of input distribution, we must actively enforce tree balance. In \textbf{Chapter~13 (Red-Black Trees)}, we introduce node coloring invariants and rotations, bounding tree height strictly by $2\log_2(n+1)$ under arbitrary dynamic updates.
\end{historybox}

\subsection{Why we randomize}

We saw in the deletion section's trap that inserting keys in sorted order yields a
height-$(n-1)$ disaster. The quicksort connection explains \emph{why}: feeding sorted keys to
BST insertion is exactly feeding quicksort the worst possible pivots (always the smallest
remaining element), the case that makes quicksort take $\Theta(n^2)$ time.

So we do what we did for quicksort: we randomize. A \emph{randomly built binary search tree}
on $n$ distinct keys is the tree you get by starting empty and inserting the keys in an order
chosen uniformly at random from all $n!$ orderings. We now ask the natural question: how tall
should we \emph{expect} such a tree to be?

\section{The expected height of a random BST}
\label{ch12:sec-expected-height}

\subsection{Why depth is just counting ancestors}

The height of a tree is the depth of its deepest node, so let us understand the depth of a
single node first. The \emph{depth} of a node is the number of edges from the root down to it
-- equivalently, the number of \emph{ancestors} it has. So if we can count, on average, how
many ancestors a key ends up with, we can get a handle on the tree's shape.

One fact powers the whole analysis. Sort the keys and rename them
$1,2,\dots,n$ by rank (so key $1$ is the smallest, key $n$ the largest). Consider two keys
$i$ and $j$ with $i < j$.

\begin{lemma}[Ancestor criterion]
\label{ch12:ancestor-criterion}
In a BST built by insertions, keys $i$ and $j$ end up in an ancestor relationship (one is an
ancestor of the other) if and only if, among all the keys with rank between $i$ and $j$
inclusive -- that is, $i, i{+}1, \dots, j$ -- the \emph{first one inserted} is $i$ or $j$
itself.
\end{lemma}

\begin{intuitionbox}
Think about who gets to be the ``pivot'' separating $i$ and $j$. Consider the block of keys
$i, i{+}1, \dots, j$. The first of these to be inserted becomes the highest of them in the
tree; it is the pivot that all the others in the block fall around. If that first-inserted key
is some $k$ strictly between $i$ and $j$, then $i$ slides into $k$'s left subtree and $j$ into
$k$'s right subtree -- they are separated forever and neither can be the other's ancestor. But
if the first-inserted key of the block is $i$ or $j$ itself, then no key ever comes between
them to split them, and one of them ends up as an ancestor of the other. Only the keys
\emph{between} $i$ and $j$ can separate them, because any key outside the block is too small
or too large to land between them in any comparison.
\end{intuitionbox}

\subsection{Counting with an indicator variable}

The block $\{i, i+1, \dots, j\}$ contains exactly $j - i + 1$ keys. Because the insertion order
is uniformly random, every one of these keys is equally likely to be the first of the block to
be inserted. Two of them ($i$ and $j$) are ``good.'' So
\[
\Pr[\,i \text{ and } j \text{ are in an ancestor relationship}\,]
   \;=\; \frac{2}{j - i + 1}.
\]
Let us turn this into a count of ancestors. Fix a key $j$ and define, for each $i \ne j$, the
indicator
\[
A_{ij} = \begin{cases} 1 & \text{if key } i \text{ is an ancestor of key } j,\\ 0 & \text{otherwise.}\end{cases}
\]
The depth of key $j$ -- its number of ancestors -- is exactly $\sum_{i \ne j} A_{ij}$. By
linearity of expectation, the \emph{expected} depth of key $j$ is $\sum_{i \ne j}
\Pr[A_{ij}=1]$. For any pair, at most one of ``$i$ ancestor of $j$'' or ``$j$ ancestor of
$i$'' can hold, and by symmetry each is equally likely; so the probability that $i$ is
specifically the ancestor of $j$ is half of $\tfrac{2}{|i-j|+1}$, namely $\tfrac{1}{|i-j|+1}$.

\subsection{The harmonic numbers appear}

Now sum over all $i \ne j$ to get the expected depth of key $j$. The keys ranked at
distance $d$ from $j$ each contribute $\tfrac{1}{d+1}$, and there are at most two such keys
(one of smaller rank, $j-d$, and one of larger rank, $j+d$), with $d$ ranging from $1$ up to
at most $n-1$. Splitting the sum into the smaller-ranked ancestors and the larger-ranked
ancestors, each half is bounded by a harmonic sum:
\[
\mathbb{E}[\text{depth of } j]
   \;=\; \sum_{i \ne j} \frac{1}{|i-j|+1}
   \;\le\; 2\sum_{d=1}^{n-1} \frac{1}{d+1}
   \;<\; 2\sum_{k=1}^{n} \frac{1}{k} \;=\; 2H_n,
\]
where $H_n$ is the \emph{$n$-th harmonic number},
\[
H_n \;=\; \sum_{k=1}^{n} \frac{1}{k} \;=\; 1 + \tfrac12 + \tfrac13 + \cdots + \tfrac1n.
\]
The harmonic number is the fingerprint of ``randomly chosen pivots,'' and it is the same
quantity that governs quicksort's expected comparison count. And $H_n$ is famously close to a
logarithm:
\[
\ln n \;<\; H_n \;\le\; \ln n + 1, \qquad\text{so}\qquad H_n = \Theta(\log n).
\]

\begin{focusbox}
Each pairwise probability $\tfrac{2}{j-i+1}$ is a piece of a harmonic sum, and harmonic sums
grow like $\ln n$. That is the entire reason a random BST is short: the same $H_n \approx \ln
n$ that makes quicksort fast makes the random tree's nodes shallow on average.
\end{focusbox}

So the expected depth of every individual node is $O(\log n)$, and summing over all $n$ nodes,
the expected total path length of a random BST is $\Theta(n \log n)$; the \emph{average} node
sits at depth $\Theta(\log n)$. But the \emph{height} is the depth of the single deepest node,
which could in principle be larger than the average. Controlling it needs one more idea.

\begin{theorem}[Expected height of a random BST]
\label{ch12:expected-height}
The expected height of a randomly built binary search tree on $n$ distinct keys is
$O(\log n)$. More precisely, as $n \to \infty$ it is asymptotic to $\alpha \ln n$ with
$\alpha \approx 4.311$ (equivalently about $2.99\,\log_2 n$), and in particular it is
$\Theta(\log n)$.
\end{theorem}

The proof in CLRS uses a clever trick: instead of analyzing the height $h$ directly, it analyzes
the \emph{exponential height} $Y_n = 2^{h}$, because expectations of products behave nicely
under the random split into left and right subtrees. One sets up a recurrence for
$\mathbb{E}[Y_n]$ -- the root falls at a uniformly random rank, splitting the keys into a left
subtree and a right subtree -- solves it to get $\mathbb{E}[Y_n] = O(n^3)$, and then uses
Jensen's inequality to pull the conclusion back down to $h$ itself. We will not grind through
every line here; the point we want you to carry is the \emph{shape} of the argument and the
inevitable appearance of $\log n$.

\begin{formalbox}
\begin{lemma}[The convexity step]
\label{ch12:jensen}
Because $f(x) = 2^x$ is convex, for any random variable $h$ we have
$2^{\,\mathbb{E}[h]} \le \mathbb{E}\!\left[2^{h}\right]$ (this is Jensen's inequality). So if we
can bound the expected exponential height $\mathbb{E}[2^h]$ by a polynomial $cn^3$, then
$2^{\mathbb{E}[h]} \le cn^3$, and taking logarithms,
\[
\mathbb{E}[h] \;\le\; \log_2(cn^3) \;=\; 3\log_2 n + O(1) \;=\; O(\log n). \qedhere
\]
\end{lemma}
\end{formalbox}

\begin{aestheticbox}
Theorem~\ref{ch12:expected-height} is genuinely good news for practice, but with a sharp
caveat. ``Expected $O(\log n)$'' assumes the insertion order is random. Real data is rarely
random -- it arrives sorted by timestamp, alphabetized, or clustered -- and an adversary who
feeds you sorted input gets the $\Theta(n)$ disaster every time. That is why production systems
do not bet on randomness; they use \emph{balanced} trees (red-black trees, AVL trees, B-trees)
that \emph{guarantee} $O(\log n)$ height no matter the input. The next chapter builds the most
famous of these. Think of the random-BST result as the reason balanced trees are worth the
trouble: it shows $O(\log n)$ is achievable, and balancing makes it dependable.
\end{aestheticbox}

\begin{walkbox}
Here is a question to mull on a walk. We proved the random tree is short \emph{on average}. But
``on average'' hides the spread. Could a random BST occasionally be tall -- say height
$\Theta(\sqrt n)$ -- with non-negligible probability, even though the average is $\Theta(\log
n)$? If you inserted a million random keys, would you ever expect to see a path of length, say,
$200$? (Hint: the expected height is about $4.31 \ln(10^6) \approx 60$, while the average
\emph{depth} is only about $\ln(10^6) \approx 14$. How far above the mean do random heights
stray?) What
would you have to prove to sleep soundly using an unbalanced BST on random data?
\end{walkbox}

\section{Putting it together: a worked lab}

\begin{labbox}
Work this entire scenario with pencil and paper before reading the answers, then check yourself.

Start with an empty BST and insert, in order: $\;50,\,30,\,70,\,20,\,40,\,60,\,80,\,35$.
\begin{enumerate}
  \item Draw the tree. (You should get $50$ at the root with two full levels below it, and $35$
        dangling as the left child of $40$.)
  \item Write the in-order walk. (Answer: $20,30,35,40,50,60,70,80$ -- check it is sorted.)
  \item Find the successor of $40$ by the algorithm. ($40$ has no right child, so climb: $40$ is
        the right child of $30$, climb to $30$; $30$ is the left child of $50$, stop -- successor
        is $50$.)
  \item Find the successor of $50$. ($50$ has a right child, so take the minimum of the right
        subtree rooted at $70$: that is $60$.)
  \item Delete $30$ (a two-child node). Its successor is $35$ (minimum of the subtree rooted at
        $40$). Splice $35$ into $30$'s place; $35$ inherits $30$'s children $20$ and $40$. Redraw
        and confirm the in-order walk is now $20,35,40,50,60,70,80$.
  \item This insertion order gave a short tree. Now imagine inserting the \emph{same eight keys}
        in sorted order $20,30,35,40,50,60,70,80$. What height do you get, and why? (A
        right-leaning chain of height $7$ -- the worst case.)
\end{enumerate}
\end{labbox}

\section*{Chapter summary}

A binary search tree stores keys so that, at every node, the left subtree holds smaller keys
and the right subtree holds larger ones. That single invariant makes an in-order walk emit the
keys in sorted order, and it makes search, minimum, maximum, successor, predecessor, insertion,
and deletion all run by tracing a single root-to-leaf (or leaf-to-root) path -- hence all cost
$O(h)$, the height. Insertion hangs a new node where a search would have failed; deletion has
three cases, the hardest replacing a doomed two-child node with its successor. The shape, and
thus the height, depends entirely on insertion order: sorted input degenerates to a height-$n$
chain. Building the tree from a random insertion order is identical to running quicksort with
those keys as pivots, and that correspondence -- through harmonic numbers and a convexity
argument -- yields an expected height of $O(\log n)$. Because real input is not random, the next
chapter trades expectation for a guarantee with balanced trees.

\section{Exercises}

\begin{exercise}[Warmup: build and walk]
Insert the keys $15, 6, 18, 3, 7, 17, 20, 2, 4, 13, 9$ into an initially empty BST in the given
order. Draw the resulting tree and write out its in-order walk. Confirm that the walk is sorted.
\end{exercise}

\begin{exercise}[Search paths]
Using your tree from the previous exercise, list the exact sequence of nodes visited by
\textsc{Tree-Search} when searching for (a) $13$, (b) $9$, and (c) $16$ (which is absent). How
many comparisons does each search make?
\end{exercise}

\begin{exercise}[Successor and predecessor by hand]
In the same tree, find the successor of $7$ and the predecessor of $13$ using the algorithms,
showing each climb or descent. Verify your answers against the in-order walk.
\end{exercise}

\begin{exercise}[Deletion cases]
Starting from your tree, delete $6$ (a two-child node), then delete $2$ (a leaf), then delete
$18$. Redraw after each deletion and confirm the in-order walk stays sorted throughout.
\end{exercise}

\begin{exercise}[Why local checking fails]
Give a binary tree with at least five nodes in which every node's left child is smaller and
every node's right child is larger than the node itself, yet the tree is \emph{not} a valid BST.
Explain precisely which node violates Definition~\ref{ch12:bst-property} and why the local check
missed it.
\end{exercise}

\begin{exercise}[Counting tree shapes]
How many \emph{structurally different} BSTs can hold the three keys $\{1,2,3\}$? Draw them all.
Then argue that the count for $n$ keys is the $n$-th Catalan number $C_n =
\frac{1}{n+1}\binom{2n}{n}$. (Hint: which key is the root, and how does that split the rest?)
\end{exercise}

\begin{exercise}[The pivot correspondence]
Take the keys $4,1,3,2,5$ and insert them in that order into a BST. Separately, run quicksort on
the array $[4,1,3,2,5]$ choosing the \emph{first} element of each subarray as the pivot. Show
that the set of pairwise comparisons made by quicksort equals the set of comparisons made by BST
insertion, pair for pair.
\end{exercise}

\begin{exercise}[Expected depth of the smallest key]
Using the ancestor criterion (Lemma~\ref{ch12:ancestor-criterion}), show that in a randomly
built BST the expected number of ancestors of the smallest key (rank $1$) is exactly
$H_n - 1 = \sum_{k=2}^{n}\frac1k$. (Hint: key $1$ has key $j$ as an ancestor exactly when $j$ is
the first inserted among ranks $1,\dots,j$.) Conclude that the expected depth of the minimum is
$\Theta(\log n)$.
\end{exercise}

\begin{exercise}[Sorting via a BST]
Describe an algorithm that sorts $n$ numbers by inserting them all into a BST and then doing an
in-order walk. What is its worst-case running time, and on what kind of input is that worst case
achieved? What is its expected running time on a random permutation, and which classic sorting
algorithm does this expected cost match -- and why do those two recurrences match?
\end{exercise}

\begin{exercise}[Challenge: in-order without recursion or parent pointers]
The \textsc{Tree-Successor} algorithm needs parent pointers. Suppose your nodes have only
$left$ and $right$ (no $parent$). Design an in-order traversal that uses an explicit stack and
prints all keys in sorted order in $\Theta(n)$ time and $O(h)$ extra space. Trace it on the
seven-key tree from the chapter opening.
\end{exercise}

\begin{exercise}[Challenge: tightening the height intuition]
We claimed the expected height is about $4.31\ln n$, yet the average \emph{depth} is about
$\ln n$ (roughly $H_n$). Explain in words why the \emph{maximum} depth (the height) should be a
constant factor larger than the \emph{average} depth, and why that constant is finite rather
than growing with $n$. (You do not need the full proof; argue from the fact that very deep paths
require an unlucky chain of pivot choices, each independently unlikely.)
\end{exercise}
